\section{Discussion}
\label{sec:discussion}

\subsection{What is fixed by the present theory}

At the level stated here, the theory fixes five structural commitments:
(i) a 4D cubic brane lattice embedded with codimension zero in a symmetric 4D Euclidean ambient, with the timelike direction selected by the Lorentzian sign structure of the brane action (\Cref{subsec:ontology-kinematics}),
(ii) a geometric amplitude--lateral coupling generated by Euclidean link lengths / the induced metric, applied symmetrically across the four ambient directions,
(iii) per-wavepacket band isolation whose adiabatically transported eigenbundle carries Berry/Wilczek--Zee connection data,
(iv) the interpretation of particles as localized standing-wave modes rather than point objects, and
(v) a retrocausal worldtube reading (\Cref{sec:bell-interpretation}) in which causality is selected by soliton chirality and entanglement is the continuity of a single 4D world-volume that branches at a V-vertex in the past.
These commitments do not yet close the theory, but they sharply constrain what kinds of completions are still allowed.

\subsection{Why the mathematical bridges must remain explicit}

The present framework only becomes interesting if its bridges are written explicitly enough to be attacked in actual
calculation.
For that reason, the continuum action, the lattice substrate, the contraction channel, the per-wavepacket band-isolation
framing, the holonomy interpretation, and the localized-mode program should not be treated as optional stylistic
embellishments.
They are the working architecture of the theory.
Removing such bridges would make the paper shorter, but it would also make the theory less testable and less coherent.

\subsection{Microscopic cubicity and macroscopic isotropy}
\label{subsec:discussion-lattice}

The dual-observer mechanism by which operational isotropy can emerge despite a cubic UV substrate is stated
canonically in \Cref{para:dual-observer}.
The discussion-level point is that this is not guaranteed.
If different branches, wavelengths, or polarization sectors feel the lattice differently, residual anisotropy and
birefringence survive as observable signatures.
The theory therefore lives or dies on a quantitative statement: whether one can identify a dominant observer sector in
which anisotropy is sufficiently suppressed and sufficiently universal, with residual birefringence as the failure
signature.

\subsection{Axis-aligned triplet phases and non-Abelian transport}
\label{subsec:discussion-axis-triplet}

The axis-aligned triplet idea is best understood as a basis adapted to the cubic substrate.
On that substrate it is natural to organize a band-isolated wavepacket triplet
$\Psi=(\psi_x,\psi_y,\psi_z)^\top$ aligned with the lattice axes, with each component arising from per-wavepacket
complex-envelope demodulation of the corresponding lateral channel.
If inter-axis mixing is negligible, each component carries an approximately independent rank-1 Berry phase.
For any nonzero mixing, however, the physical eigenmodes are mixtures of axis labels and the invariant content is the
Wilczek--Zee holonomy of the transported subspace, not a separate gauge-invariant phase ``of an axis.''
This distinction (independent-$U(1)^3$ vs Wilczek--Zee $U(3)$) is what makes the triplet sector non-trivially
different from three decoupled EM channels, and it is the theoretical reason the baryon sector is the decisive
numerical target (see \Cref{subsec:baryon-next-target}).

\subsection{Contact with lattice gauge ideas}
\label{subsec:discussion-lqcd}

Standard lattice QCD uses the lattice primarily as a regulator for a continuum gauge theory.
Here the interpretive direction is reversed: the lattice is ontic, while gauge structure is read out as connection data
of band-isolated wavepacket eigenbundles.
Despite this difference, much of the diagnostic machinery transfers directly.
Berry connections and curvatures can be computed on a discretized Brillouin zone using link variables and plaquettes,
allowing gauge-invariant comparison of Abelian and non-Abelian transport sectors.
The analogy is therefore methodological rather than ontological.

\subsection{Related emergent-$U(1)$ constructions on 3D cubic substrates}
\label{subsec:discussion-related-emergentEM}

It is worth situating the present construction against peer-reviewed precedents in which a 3D cubic-lattice
substrate is known to host an emergent $U(1)$ gauge sector at low energies.

\citet{HanKivelson2024EmergentGauge} construct a 3D cubic lattice with bond-centered atomic displacements
coupled to electrons and prove that strong electron--phonon coupling drives the system into an ice-rule
constrained ``phononic resonating-valence-bond'' phase whose effective theory near the Rokhsar--Kivelson point
is a $U(1)$ gauge theory of the form
$\mathcal{H}=\tfrac{1}{2\tilde\varepsilon}|\bm{d}|^2
+\tfrac{1}{2\tilde\mu}\big(\kappa^2|\nabla\!\cdot\!\bm{d}|^2+|\nabla\!\times\!\bm{a}|^2\big)$,
with gapless transverse-optical phonons playing the role of the emergent photon and deconfined irrationally
charged excitations playing the role of matter.
In substrate class their model is the closest published analogue of the present brane lattice: both live on a
3D cubic geometry of coupled oscillators and both produce a gapless $U(1)$ sector at low energies.
The mechanisms, however, are distinct.
There, gauge invariance is enforced as a low-energy consequence of an ice-rule constraint on bond-centered
displacements generated by an electronic sector; here, the substrate is purely mechanical and gauge structure
appears as the Berry/Wilczek--Zee holonomy of an isolated polarization band along the carrier eigenbundle.
The shared conclusion is the qualitative plausibility of an emergent gapless $U(1)$ in 3D cubic oscillator
networks; the divergence is that the present theory extends the same logic to a near-degenerate triplet whose
non-Abelian Wilczek--Zee transport realizes $U(1)\times SU(3)$ rather than $U(1)$ alone
(cf.\ \Cref{subsec:axis-triplet}), with the prestretch parameter $\alpha$ running the crossover.

A second precedent is the $U(1)$ quantum spin liquid expected in 3D pyrochlore spin ice
\citep{GingrasMcClarty2014QSI}.
There a local 2-in/2-out ice-rule constraint at each tetrahedral vertex produces an emergent compact $U(1)$
gauge theory with a gapless ``photon'' mode and deconfined magnetic monopole excitations.
The mechanism is again constraint-driven and operates on spin rather than on lattice displacement, but it
reinforces the same structural lesson: 3D substrates with a vertex-local conservation law are a natural setting
for emergent electromagnetism.

Non-Abelian gauge structure, by contrast, has been realized in \emph{classical wave} media outside the
condensed-matter setting.
\citet{Chen2019NonAbelianOptics} obtain a synthetic $SU(2)$ gauge field for optical waves directly from the
anisotropy of a homogeneous medium, and \citet{Kumar2026NonabelianElastic} design continuous elastic waveguides
whose holonomy---acquired by \emph{crossing} band degeneracies rather than by maintaining them as in the
Wilczek--Zee construction used here---transfers excitation between rods in a concatenation-order-dependent
(hence non-commuting) way.
These establish that non-Abelian holonomy on classical elastic/optical waves is physically realizable, which
lends qualitative support to the triplet-sector claim of \Cref{subsec:axis-triplet}.
The distinction is one of provenance: in both works the gauge response is \emph{engineered} into the medium by a
spatially designed anisotropy or coupling profile and demonstrated as a mode-transport effect, whereas here the
non-Abelian sector is meant to arise spontaneously from a single ontic substrate that must \emph{also} host the
$U(1)$ sector, the soliton matter sector, and the gravity-like channel together.

What none of these constructions provide is (i) non-Abelian internal structure arising spontaneously from the
same unengineered substrate that hosts the matter and gravity sectors, (ii) a Lorentz-respecting wide-soliton
sector decoupled from a Peierls--Nabarro lattice scale, or (iii) a gravity-like channel sourced by transverse
excitation through the induced metric.
These are precisely the points where the present theory takes on commitments that are absent in
\citet{HanKivelson2024EmergentGauge}, \citet{GingrasMcClarty2014QSI}, \citet{Chen2019NonAbelianOptics}, and
\citet{Kumar2026NonabelianElastic}, and that must therefore be tested on their own merits.

\subsection{Gravity as contraction from transverse excitation}

The full mechanism --- Pythagorean spring-length effect, $g_{ij}=h_\star^2\,\delta_{ij}+\partial_i u\,\partial_j u$,
and the resulting inward lateral pull under homogeneous prestretch --- is derived in \Cref{subsec:gravity-channel}.

What remains open is the quantitative reduction.
One would like to derive a Newtonian limit in which an effective scalar potential built from coarse-grained strain or
energy density obeys a Poisson-type equation with a universal coupling constant, and then relate that constant to $G$.
On the simulation side, the central check is straightforward: create a localized transverse excitation, measure the
induced steady contraction field, and quantify its action on long-wavelength probe packets.

\subsection{Quantum phenomenology without fundamental probability}

By construction, the underlying dynamics is deterministic.
The probabilistic character of quantum measurements is attributed to the interaction of localized modes with wave packets
that are constrained by Fourier localization and by the finite Brillouin zone, together with nonlinear exchange events
in which small phase and amplitude differences can decide whether energy transfer occurs.
In this view, apparent point-like interactions reflect the structure of wave packets and nonlinear exchange,
not the existence of fundamental point particles.

\subsection{Priority open problems}

Several problems remain before the theoretical framework could be regarded as a mature replacement attempt:
\begin{itemize}
\item establishing a universal observer sector with sufficiently suppressed long-wavelength anisotropy,
\item deriving a quantitative contraction-channel reduction and matching it to an effective gravitational constant,
\item identifying parameter regimes that simultaneously support band-isolated adiabatic wavepackets
and robust localized modes,
\item determining whether the cubic triplet sector really supports stable baryon-like confined modes,
\item clarifying how far-field Abelian holonomy and internal non-Abelian transport map onto observed particle labels,
\item understanding multi-particle states and interaction processes in the deterministic substrate.
\end{itemize}
The paper is now best read as a sharpened core theory plus a numerical research program whose first decisive test is the
existence or failure of stable triplet-confined baryon-like modes.
